\section{Introducción}

\subsection{I. Descripción de la empresa y/o Área afectada por el proyecto a implementar}

\subsubsection{Nombre de la empresa y/o Área}
\textbf{Empresa:} Global GIS Service (GGS) \\
\textbf{Área afectada:} Investigación y Desarrollo (I+D)

\subsubsection{Misión, Visión de la empresa}
\textbf{Misión:} Desarrollar soluciones geoespaciales innovadoras que optimicen la toma de decisiones estratégicas en organizaciones públicas y privadas, integrando tecnología GIS con metodologías ágiles y sostenibles.

\textbf{Visión:} Ser líderes en Latinoamérica en el desarrollo de plataformas GIS interoperables, escalables y centradas en el usuario, contribuyendo a la transformación digital del sector educativo, gubernamental y empresarial.

\subsubsection{Objetivos Estratégicos de la empresa y/o Área}
\textbf{Objetivos de GGS:}
\begin{itemize}
  \item Expandir la adopción de soluciones GIS en sectores clave como educación, gobierno y servicios públicos.
  \item Garantizar la interoperabilidad entre plataformas GIS y sistemas de terceros.
  \item Fortalecer la relación con clientes mediante procesos de validación colaborativos y comunicación efectiva.
  \item Promover la capacitación continua de usuarios finales para asegurar el uso eficiente de las soluciones implementadas.
\end{itemize}

\textbf{Objetivos del Área de I+D:}
\begin{itemize}
  \item Diseñar una metodología de trabajo sostenible y replicable para el desarrollo de proyectos GIS.
  \item Fomentar la innovación mediante la incorporación de tecnologías emergentes.
  \item Asegurar la calidad técnica de los entregables a través de revisiones periódicas y pruebas de integración.
  \item Generar conocimiento interno que permita escalar soluciones de manera eficiente en nuevos mercados.
\end{itemize}

\subsubsection{Objetivos del proyecto}
\begin{itemize}
  \item Implementar una solución GIS interoperable que integre datos del cliente en tiempo real para la toma de decisiones.
  \item Capacitar a usuarios finales en el uso de herramientas WebGIS para consultas y análisis espaciales básicos.
  \item Entregar un panel de control con métricas de uso y cumplimiento de actividades obligatorias.
  \item Documentar arquitectura, procesos y manuales de usuario para operación y mantenimiento.
\end{itemize}

\subsubsection{Alineamiento del proyecto con los objetivos de la empresa}
El proyecto ``Consultoría Integral en Soluciones GIS'' se alinea directamente con los objetivos estratégicos de GGS y el área de I+D, ya que:
\begin{itemize}
  \item Promueve la interoperabilidad tecnológica mediante la integración de plataformas GIS con sistemas del cliente.
  \item Establece una metodología sostenible y replicable a través de sesiones prácticas, manuales claros y validaciones técnicas.
  \item Fomenta la gestión del conocimiento mediante reuniones técnicas mensuales e informes especializados.
  \item Impulsa la transformación digital del cliente mediante la automatización de procesos de validación y comunicación.
\end{itemize}

\subsection{II. Aplicación: Proyecto}

\subsubsection{Descripción del alcance del proyecto}

\paragraph{(i) Alcance general del Proyecto a Implementar}
El proyecto consiste en el diseño e implementación de servicios de consultoría especializados en Sistemas de Información Geográfica (GIS), orientados a mejorar la toma de decisiones mediante análisis espacial, interoperabilidad tecnológica y capacitación de usuarios. Incluye:
\begin{itemize}
  \item Análisis de necesidades y requisitos del cliente.
  \item Desarrollo de soluciones GIS personalizadas.
  \item Integración de plataformas GIS con sistemas existentes.
  \item Capacitación básica a usuarios finales.
  \item Documentación técnica y manuales de usuario.
\end{itemize}

\paragraph{(ii) Historias de usuario del proyecto}
\begin{itemize}
  \item HU001: Como cliente, quiero integrar mi plataforma de gestión con el sistema GIS para visualizar datos geoespaciales en tiempo real. \\
  \textit{Criterios de aceptación:} (a) Conectores configurados a al menos una fuente de datos; (b) Visualización en mapa y tabla en menos de 3 segundos de latencia; (c) Registro de eventos de integración.
  \item HU002: Como usuario final, necesito capacitación básica en el uso de WebGIS para realizar consultas espaciales autónomas. \\
  \textit{Criterios de aceptación:} (a) Módulo de capacitación entregado; (b) 80\% de participantes aprueban evaluación; (c) Guía rápida disponible.
  \item HU003: Como administrador, requiero un panel de control con métricas de uso y cumplimiento de actividades obligatorias en la plataforma GIS. \\
  \textit{Criterios de aceptación:} (a) Dashboard con KPIs definidos; (b) Exportación de reportes; (c) Control de acceso por roles.
\end{itemize}

\paragraph{(iii) Roles Scrum}
\begin{itemize}
  \item \textbf{Product Owner:} Representante del cliente o área solicitante (ej: Dra. Rosmery R.). Responsable de maximizar el valor del producto, priorizar el backlog y aceptar entregables.
  \item \textbf{Scrum Master:} Facilitador del equipo, responsable de eliminar impedimentos, asegurar el marco Scrum y mejorar el flujo de trabajo (ej: Luis Huatay).
  \item \textbf{Development Team:} Especialistas GIS, analistas técnicos, desarrolladores. Responsables de construir el incremento, pruebas, documentación y despliegues.
\end{itemize}

\subsubsection{Artefactos Scrum}

\paragraph{(i) Product Backlog}
Lista priorizada de funcionalidades y requerimientos del proyecto, derivada del alcance inicial. Ejemplos de ítems:
\begin{itemize}
  \item Diseño de arquitectura de base de datos GIS.
  \item Desarrollo de mapas temáticos y dashboards.
  \item Implementación de integración con sistemas del cliente.
  \item Redacción de manuales técnicos y de usuario.
\end{itemize}

\paragraph{(ii) Sprint Backlog}
Conjunto de tareas seleccionadas del Product Backlog para un sprint específico. Ejemplo para el Sprint 1:
\begin{itemize}
  \item T1: Inventariar sistemas del cliente para integración.
  \item T2: Diseñar prototipo de dashboard GIS.
  \item T3: Validar interoperabilidad con plataformas existentes.
\end{itemize}

\paragraph{(iii) Increment}
Producto funcional entregado al final de cada sprint. Ejemplo:
\begin{itemize}
  \item Sprint 1: Prototipo de dashboard GIS básico, con integración inicial a una fuente de datos.
\end{itemize}

\subsubsection{Descripción de eventos Scrum}

\paragraph{(i) Sprint Planning}
\textbf{Objetivo:} Planificar el trabajo del sprint, seleccionando historias del Product Backlog. \\
\textbf{Duración:} 2 horas. \\
\textbf{Participantes:} Product Owner, Scrum Master, Equipo de Desarrollo. \\
\textbf{Resultado:} Objetivo del sprint, Sprint Backlog, definición de Done aplicable.

\paragraph{(ii) Daily (detalle de un Daily)}
\textbf{Fecha:} 12 de noviembre de 2025, 10:00 AM. \\
\textbf{Duración:} 15 minutos. \\
\textbf{Reportes:}
\begin{itemize}
  \item Ana: Ayer: Revisé requisitos de integración. Hoy: Diseñaré flujo de datos. Impedimentos: Ninguno.
  \item Carlos: Ayer: Avancé en desarrollo de API. Hoy: Terminaré pruebas unitarias. Impedimentos: Esperando acceso a servidor de pruebas.
  \item Luis: Ayer: Configuré CI/CD para despliegues del dashboard. Hoy: Integraré métricas. Impedimentos: Falta token de acceso.
\end{itemize}

\paragraph{(iii) Review (cómo se llevará a cabo)}
\textbf{Objetivo:} Presentar el incremento al cliente y recibir feedback. \\
\textbf{Participantes:} Equipo Scrum, cliente, stakeholders clave. \\
\textbf{Actividades:} Demostración en vivo, discusión de mejoras, ajustes al backlog.

\paragraph{(iv) Retrospective (de uno de los Sprint)}
\textbf{Sprint 1 – Técnica:} Start, Stop, Continue. \\
\textbf{Resultados:}
\begin{itemize}
  \item Start: Incluir pruebas de accesibilidad en el pipeline de QA.
  \item Stop: Implementar cambios sin aprobación previa del PO.
  \item Continue: Mantener Daily Scrums a las 10:00 AM.
\end{itemize}

\paragraph{(v) Product Backlog Refinement}
\textbf{Objetivo:} Ajustar y priorizar ítems del backlog antes de cada sprint planning. \\
\textbf{Frecuencia:} Semanal. \\
\textbf{Responsable:} Product Owner con apoyo del equipo.

\subsubsection{Tablero Kanban del proyecto: Indicar todas las tareas y sus estados}
\begin{center}
\begin{tabular}{lllc}
\hline
\textbf{Tarea ID} & \textbf{Descripción} & \textbf{Estado} & \textbf{Responsable} \\
\hline
T1 & Análisis de requisitos del cliente & Done & Ana \\
T2 & Diseño de arquitectura GIS & In Progress & Carlos \\
T3 & Desarrollo de API de integración & In Progress & Luis \\
T4 & Pruebas de interoperabilidad & To Do & María \\
T5 & Redacción de manual técnico & To Do & Pedro \\
T6 & Capacitación a usuarios finales & To Do & Scrum Master \\
\hline
\end{tabular}
\end{center}
\noindent Nota: El tablero refleja el 100\% de las tareas planificadas para el proyecto en este ciclo, con su estado y responsable actual.

\subsection{III. Conclusiones}
El proyecto ``Consultoría Integral en Soluciones GIS'' representa una oportunidad estratégica para alinear los objetivos de Global GIS Service con las necesidades del mercado, mediante la implementación de una solución interoperable, escalable y centrada en el usuario. La aplicación del marco Scrum permitirá gestionar el desarrollo de manera iterativa e incremental, asegurando entregas tempranas, retroalimentación continua y adaptabilidad ante cambios. La combinación de una planificación robusta (PMI) con ejecución ágil (Scrum) garantizará no solo el cumplimiento de los plazos y presupuestos, sino también la satisfacción del cliente y el fortalecimiento del posicionamiento de GGS en el sector geoespacial latinoamericano.